\documentclass[11pt]{article}

\usepackage[a4paper,margin=1in]{geometry}
\usepackage{amsmath,amssymb,mathtools}
\usepackage[hidelinks]{hyperref}
\usepackage{enumitem}
\usepackage{xcolor}
\usepackage{listings}

\lstset{
  basicstyle=\ttfamily\small,
  breaklines=true,
  columns=fullflexible,
  frame=single,
  xleftmargin=0.5em,
  xrightmargin=0.5em
}

\title{Precise Feedback on the Improved Version\\
\large Especially on the Irrational-Slope Section}
\author{}
\date{\today}

\begin{document}
\maketitle

\section*{Overall assessment}

The improved version is substantially stronger than the previous one.  In
particular, the new Section ``The Irrational Case'' now uses the correct
cut-and-project mechanism:
\[
    \text{physical injectivity}
    \;+
    \text{internal density}
    \;+
    \text{absence of non-trivial interval stabilizers}.
\]
The previous unjustified lifting step has essentially been repaired.

The rational part of the paper still appears mathematically sound.  The
period formulas, the residue reduction, the generic minimality argument, and
the set-valued comparison remain coherent.  I do not see a fatal flaw in the
main rational theorem.

That said, I would not yet send the paper to a specialist without a few
precise corrections.  The most concrete mathematical problem is in the
figures: an accepted lattice point is missing.  The irrational proof is now
basically correct, but one implication should be written out more explicitly.

\section{Must-fix point: the figures omit an accepted lattice point}

In both TikZ figures the chosen example is
\[
    \alpha=2,\qquad \beta=1,\qquad \omega=1.
\]
The strip condition is therefore
\[
    2(x-1) \le y \le 2x+1.
\]
For \(x=0\), this gives
\[
    -2 \le y \le 1.
\]
Your visible lattice includes \(y=-1\), and the point \((0,-1)\) lies inside
the strip.  It is therefore an accepted lattice point.  However, Figure~1
currently lists only
\begin{lstlisting}
0/0, 0/1
\end{lstlisting}
for \(x=0\), omitting
\begin{lstlisting}
0/-1
\end{lstlisting}

This is not merely cosmetic: a specialist could notice it immediately.

\subsection*{Correction for Figure 1}

Add \texttt{0/-1} to both loops, the blue accepted-point loop and the
projection loop.  Replace the accepted-point list by:

\begin{lstlisting}
\foreach \x/\y in {0/-1, 0/0, 0/1, 1/0, 1/1, 1/2, 1/3,
                   2/2, 2/3, 2/4, 2/5,
                   3/4, 3/5, 3/6, 3/7} {
  \fill[blue!65!black] (\x,\y) circle (2.6pt);
}
\end{lstlisting}

and likewise replace the projection loop by:

\begin{lstlisting}
\foreach \x/\y in {0/-1, 0/0, 0/1, 1/0, 1/1, 1/2, 1/3,
                   2/2, 2/3, 2/4, 2/5,
                   3/4, 3/5, 3/6, 3/7} {
  \pgfmathsetmacro{\px}{(\x + 2*\y)/5}
  \pgfmathsetmacro{\py}{2*(\x + 2*\y)/5}
  \draw[gray!65, thin] (\x,\y) -- (\px,\py);
  \fill[red!70!black] (\px,\py) circle (1.9pt);
}
\end{lstlisting}

\subsection*{Correction for Figure 2}

For the same example, the track \(r=1\) is
\[
    r=2x-y=1,
    \qquad\text{equivalently}\qquad
    y=2x-1.
\]
Thus it contains \((0,-1)\).  The orange line should pass through this point.
Replace

\begin{lstlisting}
\draw[orange!80!black, thick] (0.5,0) -- (3,5);
\end{lstlisting}

by

\begin{lstlisting}
\draw[orange!80!black, thick] (0,-1) -- (3,5);
\end{lstlisting}

and replace the orange accepted points by

\begin{lstlisting}
\foreach \x/\y in {0/-1, 1/1, 2/3, 3/5}
  \fill[orange!80!black] (\x,\y) circle (2.6pt);
\end{lstlisting}

\section{The irrational-slope proof}

The new proof of Proposition~\texttt{\textbackslash ref\{prop:irrational\}}
is now based on the right idea.  The definitions
\[
    \tilde p(x,y)=x+ay,
    \qquad
    s(x,y)=y-ax
\]
are appropriate.  The strip becomes
\[
    -a\omega \le s(x,y) \le \omega,
\]
and for irrational \(a\), the map \(\tilde p\) is injective on
\(\mathbb Z^2\).  The internal group
\[
    s(\mathbb Z^2)=\{n-am:m,n\in\mathbb Z\}
\]
is dense in \(\mathbb R\).  These are exactly the ingredients one should use.

However, Step~2 is still slightly compressed.  You write that if the projected
gap sequence has a finite period \(\lambda\), and \(\sigma\) is the sum of one
period of consecutive gaps, then
\[
    \Lambda+\sigma=\Lambda.
\]
This is true, but I recommend spelling out the enumeration argument.

\subsection*{Recommended replacement for the beginning of Step 2}

Replace the first paragraph of Step~2 by the following more explicit version:

\begin{lstlisting}
Suppose, for contradiction, that the projected gap sequence has a finite
period \(\lambda \in \mathbb{N}\).  Let
\[
    \Lambda := \tilde p(A)\subset\mathbb R.
\]
Since \(\Lambda\) is locally finite and unbounded in both directions, write
its increasing enumeration as \((\ell_i)_{i\in\mathbb Z}\).  If
\[
    g_i := \ell_{i+1}-\ell_i
\]
denotes the gap sequence and \(g_{i+\lambda}=g_i\) for all \(i\), then the
sum
\[
    \sigma := g_i+g_{i+1}+\cdots+g_{i+\lambda-1}
\]
is independent of \(i\).  Hence
\[
    \ell_{i+\lambda}-\ell_i=\sigma
    \qquad\text{for all }i\in\mathbb Z.
\]
Since all gaps are positive, \(\sigma>0\), and therefore
\[
    \Lambda+\sigma=\Lambda.
\]
\end{lstlisting}

Then continue with:

\begin{lstlisting}
Choose any \(z_0\in A\); since \(\tilde p(z_0)+\sigma\in\Lambda\), ...
\end{lstlisting}

This makes the implication from periodic gaps to translation symmetry fully
transparent.

\section{Local finiteness in the irrational case}

The phrase ``projected gap sequence'' in the irrational case presupposes that
\(\Lambda=\tilde p(A)\) admits a bi-infinite increasing enumeration.  This is
standard for regular one-dimensional cut-and-project sets, but for maximum
precision it should be mentioned.

A concise sentence before Proposition~\texttt{\textbackslash ref\{prop:irrational\}}
would suffice:

\begin{lstlisting}
For irrational \(a\), the projected set is locally finite and unbounded in
both directions, so its gap sequence is defined by the increasing
\(\mathbb Z\)-indexed enumeration.
\end{lstlisting}

If you want to be more explicit, use:

\begin{lstlisting}
The strip has bounded width in the internal direction, while the lattice is
discrete; hence bounded intervals in physical space contain only finitely
many projected points.  The set is unbounded in both directions because
\(s(\mathbb Z^2)\) is dense and the window \(W\) has nonempty interior.
\end{lstlisting}

The first sentence is probably enough for this paper.

\section{Step 3 in the irrational proof}

Step~3 is essentially correct.  The invariance
\[
    \bigl(W\cap s(\mathbb Z^2)\bigr)+\tau
    =
    W\cap s(\mathbb Z^2)
\]
combined with density of \(s(\mathbb Z^2)\) rules out every nonzero
\(\tau\).  This is the standard internal-space argument.

One small stylistic improvement: instead of

\begin{lstlisting}
pick \(t \in s(\mathbb{Z}^2) \cap W\) with \(t > \omega - \tau\)
\end{lstlisting}

write

\begin{lstlisting}
By density, choose \(t\in s(\mathbb Z^2)\cap W\) sufficiently close to
\(\omega\) so that \(t+\tau>\omega\).
\end{lstlisting}

This avoids any possible edge-case discomfort when \(\tau\) is large.  Use the
analogous wording near the lower endpoint for \(\tau<0\).

\section{The period dichotomy remark}

The dichotomy remark should explicitly include the hypothesis \(\omega>0\).
Your rational theorem allows \(\omega\ge0\), while the irrational proposition
assumes \(\omega>0\).  For \(\omega=0\) and irrational slope, the strip
degenerates to the irrational line, which contains only \((0,0)\) as a lattice
point, so the bi-infinite gap sequence is not defined in the same way.

Replace

\begin{lstlisting}
Within the family of positive slopes considered here,
Corollary~\ref{cor:finite} and Proposition~\ref{prop:irrational} together
form a complete dichotomy:
\end{lstlisting}

by

\begin{lstlisting}
Within the family of positive slopes considered here, and for \(\omega>0\),
Corollary~\ref{cor:finite} and Proposition~\ref{prop:irrational} together
form a complete dichotomy:
\end{lstlisting}

Then the equivalence
\[
    \text{finite gap period}
    \quad\Longleftrightarrow\quad
    \text{rational slope}
\]
is precise.

\section{Organization paragraph}

The introduction now mentions the paper organization, but it does not yet
mention the new irrational section.  Insert one sentence:

\begin{lstlisting}
Section~\ref{sec:corollaries} records some immediate corollaries.
Section~\ref{sec:irrational} records the complementary irrational case.
Section~\ref{sec:lean} describes the Lean~4 formalization.
\end{lstlisting}

\section{Terminology: window parameter, not window width}

Near the end of the corollaries you still write:

\begin{lstlisting}
This occurs when the window width is tuned so
that the number of tracks is an exact multiple of \(D = \alpha^2+\beta^2\).
\end{lstlisting}

Since \(\omega\) is a parameter of the defining inequalities, not literally
the orthogonal width, I recommend:

\begin{lstlisting}
This occurs when the window parameter is tuned so
that the number of tracks is an exact multiple of \(D = \alpha^2+\beta^2\).
\end{lstlisting}

\section{Lean table caption: check path capitalization}

The commit hash in the table caption is a good improvement.  However, please
check the path capitalization.  You currently write:

\begin{lstlisting}
\texttt{LEAN/CutAndProject/CutAndProject/Basic.lean}
\end{lstlisting}

Earlier versions used something like \texttt{Lean/CutAndProject/}.  GitHub
paths are case-sensitive.  If the repository uses \texttt{Lean}, the caption
should say exactly:

\begin{lstlisting}
\texttt{Lean/CutAndProject/CutAndProject/Basic.lean}
\end{lstlisting}

or whatever the actual path is.  Also, if possible, use the full commit hash
rather than only \texttt{0b12c26}.

\section{Residue-distribution wording}

This sentence is correct but slightly informal:

\begin{lstlisting}
As \(r\) runs through \(N\) consecutive integers, the residues \(mr \bmod D\)
cycle through all \(D\) classes with period~\(D\).
\end{lstlisting}

I would write instead:

\begin{lstlisting}
As \(r\) runs through consecutive integers, the residues \(mr\bmod D\)
advance by the fixed unit \(m\).  Hence every block of \(D\) consecutive
values of \(r\) hits each residue class exactly once.
\end{lstlisting}

This is more precise because the cyclic order is the \(m\)-stepping order,
not necessarily the natural order.

\section{Acknowledgments}

The acknowledgement still contains the phrase ``proof construction'':

\begin{lstlisting}
The author used ChatGPT, Gemini, and Claude as AI assistants during mathematical
exploration, drafting, proof construction, and the creation of the Lean formalization.
\end{lstlisting}

I would soften this.  A safer and more professional version is:

\begin{lstlisting}
The author used AI assistants during exploratory computations, drafting,
and the preparation of the Lean formalization.  The author has independently
verified the mathematical content and accepts full responsibility for the
results and any errors.
\end{lstlisting}

This is transparent without inviting the interpretation that the proofs were
outsourced.

\section{Minor consistency: nonuniform vs. non-uniform}

You use both ``nonuniform'' and ``non-uniform''.  This is not a mathematical
issue, but choose one spelling throughout.  I would use ``non-uniform'' in
prose and theorem captions.  The Lean names can of course retain their exact
identifier spelling.

\section*{Final verdict}

The paper is now close to being ready for external expert feedback.  The
rational part remains solid, and the irrational section now has the right
mathematical structure.  Before sending it, I would make the following changes:

\begin{enumerate}[label=\arabic*.]
    \item Add the missing accepted point \((0,-1)\) to both figures.
    \item Expand the proof that a finite gap period implies
          \(\Lambda+\sigma=\Lambda\).
    \item Mention local finiteness and bi-infinite enumerability in the
          irrational case.
    \item State the rational/irrational dichotomy for \(\omega>0\).
    \item Add Section~\texttt{\textbackslash ref\{sec:irrational\}} to the
          organization paragraph.
    \item Replace ``window width'' by ``window parameter''.
    \item Check capitalization of the Lean repository path and preferably use
          the full commit hash.
    \item Soften the AI acknowledgement.
\end{enumerate}

With these changes, I would consider the version mathematically serious and
suitable for sending to a specialist for feedback.

\end{document}
